\section{问题背景}

% Page 1: 什么是服务网络设计
\begin{frame}{服务网络设计（Service Network Design）}
\begin{block}{定义}
\justifying
服务网络设计是指在给定物理网络上，决定\hl{开设哪些服务}（路线、频率、时刻、资源分配），以最优化运营目标（利润、成本、服务水平）的规划问题。
\end{block}

\vspace{0.3cm}
\begin{block}{现实应用}
\begin{itemize}
    \item \hl{航空}：航班时刻表与机队分配
    \item \hl{城际铁路}：列车开行方案
    \item \hl{城市公交}：线路频率与车辆调度
    \item \hl{货运物流}：干线运输网络规划
\end{itemize}
\end{block}
% \begin{columns}[T,totalwidth=\textwidth]
% \column{0.48\textwidth}

% % \column{0.48\textwidth}
% % \begin{block}{共同特征}
% % \begin{itemize}
% %     \item 服务供给决定了顾客可选的出行方案
% %     \item 顾客选择反过来影响服务的盈利性
% %     \item 供给与需求之间存在\hl{反馈耦合}
% % \end{itemize}
% \end{block}
% \end{columns}
\end{frame}


% Page 2: 在SND中引入顾客选择
\begin{frame}{考虑顾客选择的服务网络设计}
\begin{columns}[T,totalwidth=\textwidth]
\column{0.45\textwidth}
在传统的服务网络设计中，顾客需求通常被视为\hl{外生输入}。
\begin{block}{考虑顾客选择的服务网络设计}
\justifying
需要同时考虑：
\begin{itemize}
    \item 开哪些服务、分配哪类资源
    \item 顾客如何选择
\end{itemize}
\hl{供给决策} 与 \hl{需求分配} 相互反馈。
\end{block}

% \vspace{0.3cm}
% \begin{block}{核心矛盾}
% \centering
% \Large
% \textbf{模型更真实}\par
% \vspace{0.15cm}
% $\Downarrow$\par
% \vspace{0.15cm}
% \Large
% \textcolor{BadRed}{\textbf{直接求解更困难}}
% \end{block}

\column{0.53\textwidth}
\begin{center}
\begin{tikzpicture}[node distance=0.5cm, every node/.style={font=\small}]
    \node[draw=structure,rounded corners,thick,fill=structure!8,minimum width=4.2cm,minimum height=1cm,align=center] (supply) {供给决策\\开哪些服务？分配什么资源？};
    \node[draw=structure,rounded corners,thick,fill=structure!8,below=1.2cm of supply,minimum width=4.2cm,minimum height=1cm,align=center] (demand) {需求分配\\顾客如何选择 itinerary？};
    \draw[-{Latex[length=2.5mm]},thick] (supply.east) to[bend left=40] node[right,font=\scriptsize]{服务方案} (demand.east);
    \draw[-{Latex[length=2.5mm]},thick] (demand.west) to[bend left=40] node[left,font=\scriptsize]{盈利结果} (supply.west);
\end{tikzpicture}
\end{center}
\end{columns}
\end{frame}


% Page 3: 为什么难
\begin{frame}{难点分析}
\begin{columns}[T,totalwidth=\textwidth]
\column{0.85\textwidth}
\begin{block}{建模方式}
\begin{itemize}
    \item 基于\hl{混合整数规划}（MIP）
    \item 网络层：多商品流（MCF）建模服务与资源
    \item 需求层：离散选择模型建模顾客行为
\end{itemize}
\end{block}

\vspace{0.2cm}
\begin{block}{计算困难的来源}
\begin{itemize}
    \item 问题本身是 \hl{NP-hard}
    \item 顾客选择是\hl{内生需求}，不是固定流量——供给变则需求变
    \item 决策变量同时包含二元服务变量 $\fleetAssignVar$ 和连续选择变量 $\shareVar$，两者耦合
    \item 可选的行程（itinerary）数量随网络规模呈现极快的增长
\end{itemize}
\end{block}

% \column{0.44\textwidth}
% \begin{center}
% \begin{tikzpicture}[node distance=0.6cm, every node/.style={font=\small}]
%     \node[draw=structure,rounded corners,thick,fill=structure!8,minimum width=4.5cm,minimum height=0.9cm,align=center] (a) {网络层：MCF + 整数变量};
%     \node[draw=structure,rounded corners,thick,fill=structure!8,below=of a,minimum width=4.5cm,minimum height=0.9cm,align=center] (b) {需求层：离散选择模型};
%     \node[draw=BadRed,rounded corners,thick,fill=BadRed!8,below=of b,minimum width=4.5cm,minimum height=1cm,align=center] (c) {耦合 $\Rightarrow$ 变量爆炸\newline LP 松弛极弱};
%     \draw[-{Latex[length=2.5mm]},thick] (a) -- (b);
%     \draw[-{Latex[length=2.5mm]},thick] (b) -- (c);
%     \draw[-{Latex[length=2.5mm]},thick,dashed,bend right=40] (c.west) to node[left,font=\scriptsize]{反馈} (a.west);
% \end{tikzpicture}
% \end{center}
\end{columns}
\end{frame}


% Page 4: 现有方法的不足
\begin{frame}{现有方法的不足}
\begin{columns}[T,totalwidth=\textwidth]
\column{0.65\textwidth}
\begin{block}{直接 MIP}
\begin{itemize}
    \item 变量太多，求解器难以直接处理
    \item 退化严重，分支定界树膨胀
\end{itemize}
\end{block}

\vspace{0.18cm}
\begin{block}{Dynamic Discretization Discovery (DDD) \cite{boland2017continuous}}
\begin{itemize}
    \item 核心思想：对\hl{时空网络做聚合}，逐步细化时间点
    \item 能有效压缩物理时空网络
    \item 但\hl{无法处理行程维度的数量爆炸}
\end{itemize}
\end{block}

% \column{0.48\textwidth}
% \begin{center}
% \begin{tikzpicture}[node distance=0.7cm]
%     \node[draw=BadRed,rounded corners,thick,fill=BadRed!8,minimum width=4.7cm,minimum height=0.9cm,align=center] (a) {Raw model 太大};
%     \node[draw=WarnOrange,rounded corners,thick,fill=WarnOrange!12,below=of a,minimum width=4.7cm,minimum height=0.9cm,align=center] (b) {DDD 只压缩时间层};
%     \node[draw=GoodGreen,rounded corners,thick,fill=GoodGreen!10,below=of b,minimum width=4.7cm,minimum height=1cm,align=center] (c) {需要同时压缩\newline 物理网络 + 顾客行程层};
%     \draw[-{Latex[length=2.5mm]},thick] (a) -- (b);
%     \draw[-{Latex[length=2.5mm]},thick] (b) -- (c);
% \end{tikzpicture}
% \end{center}
\end{columns}
\end{frame}


% Page 5: 研究目标
\begin{frame}{研究目标}
\begin{block}{本文要解决的核心问题}
\centering
\large
设计一种精确优化方法，使得：
\end{block}

\vspace{0.3cm}
\begin{columns}[T,totalwidth=\textwidth]
\column{0.32\textwidth}
\begin{block}{\centering 模型真实性}
\centering
\vspace{0.1cm}
考虑选择模型 \\[0.2cm]
供给—需求反馈耦合\\[0.2cm]
\vspace{0.1cm}
\end{block}

\column{0.32\textwidth}
\begin{block}{\centering 计算可行性}
\centering
\vspace{0.1cm}
模型规模可控\\[0.2cm]
实际可求解\\[0.2cm]
\vspace{0.1cm}
\end{block}

\column{0.32\textwidth}
\begin{block}{\centering 理论保证}
\centering
\vspace{0.1cm}
全局最优性保证\\[0.2cm]
精确方法\\[0.2cm]
\vspace{0.1cm}
\end{block}
\end{columns}
\end{frame}
